\section{Implementation Details}
\label{app:implementation-details}

\paragraph{Random seeds.}
All quantitative downstream task evaluations in both the main text and appendix are conducted across five random seeds: 0, 42, 1234, 3407, and 2024. Qualitative visualizations, t-SNE feature projections, and subset diagnostic analyses use a fixed default seed of 42 unless explicitly specified otherwise.

\paragraph{Hardware and datasets.}

Each experimental run is executed on a single NVIDIA H200 GPU; no run uses multi-GPU training. For image classification, we follow \safefix~\citep{xu2025safefix}: a ResNet-18 classifier~\citep{he2016deep} on an ImageNet subset~\citep{deng2009imagenet} with 10 classes, 1,300 images per class, and an 8:1:1 train/validation/test split, denoted ImageNet10. The ImageNet10 classes are \texttt{backpack}, \texttt{barber chair}, \texttt{coffee mug}, \texttt{desk}, \texttt{electric guitar}, \texttt{park bench}, \texttt{pitcher}, \texttt{purse}, \texttt{rocking chair}, and \texttt{water bottle}. We also evaluate CelebA~\citep{liu2015deep} lipstick-wearing classification with a ResNet-18-based binary classifier and 80,000/10,000/10,000 train/validation/test images. For pose estimation, we follow HiBug2~\citep{chen2025hibug2} on COCO person keypoints~\citep{lin2014microsoft}, using Keypoint R-CNN~\citep{mask-rcnn-2017} with a ResNet-50-FPN backbone~\citep{he2016deep}, 20,000 source training images, and 1,173-image held-out validation and test splits. For object detection, we use KITTI~\citep{geiger2012we} Car and Pedestrian instances converted to YOLO format with 5,000/1,240/1,241 train/validation/test images, and train YOLOv8n.
\paragraph{Training and repair schedule.}
For ImageNet10 classification, the diagnosis-stage ResNet-18 model is trained with cross-entropy loss and Adam~\citep{kingma2015adam}, using learning rate $1\times10^{-4}$, weight decay $1\times10^{-3}$, batch size 64, and 30 epochs. During repair, we train with the selected source subset and the generated repair images for 15 epochs using learning rate $1\times10^{-5}$. For CelebA, COCO, and KITTI, the repair-stage epoch count is set to half of the corresponding diagnosis-stage training schedule, with the same repair-stage learning rate $1\times10^{-5}$.

\paragraph{Image generation.}
For the generated repair images used in the main results, we use \controlnet~\citep{zhang2023adding} with the Stable Diffusion 1.5 checkpoint, whose backbone follows the latent diffusion architecture~\citep{rombach2022high}. The model is conditioned on soft HED boundaries~\citep{xie2015holistically} extracted from source images and uses 30 DDIM inference steps~\citep{ddim}. The soft boundary condition helps preserve object structure during edits, which is useful for attribute changes such as recoloring, material changes, and weather or style edits.

\paragraph{Attribute proposal and assignment.}
We use vision-language models (VLMs) in two separate roles. First, a VLM proposes the semantic attribute vocabulary $\attrs$ for a task. For example, in KITTI object detection, GPT-5~\citep{openai2025gpt5} is used to define candidate visual attributes such as color, weather, scale, viewpoint, occlusion, truncation, and motion state. Second, a VLM assigns attribute values to source and generated images. In the experiments reported in this paper, we use Qwen3-VL~\citep{bai2025qwen3vl} for this assignment step. The selector fields $\selattrs$ are the fields passed to the attribute-selection procedure, while the target fields $\targetattrs$ define the diagnosed conjunction and generated repair pool. The main selector fields are color, background, and texture for ImageNet10; hair color, skin, and emotion for CelebA; clothes color and weather for COCO; and color and weather for KITTI. The generated target fields are color and texture for ImageNet10, the three listed fields for CelebA, clothes color and weather for COCO, and color and weather for KITTI.

\paragraph{Selection groups for dense prediction.}
The group used by the balancing term is a selection group rather than a universal benchmark class. For ImageNet10 and CelebA, each image uses its benchmark class label as $g_i$. For COCO pose estimation, the dataset contains only the person category, so we construct a pose-complexity proxy group by counting valid person annotations in each image and assigning the counts to up to four quantile bins. This proxy label is used only for source selection; the task labels remain person keypoint annotations, and the VLM attributes are assigned at the image level. For KITTI, each image is assigned a label-signature group determined by the set of object class IDs present in its YOLO label file. A multi-object image therefore contributes once to the group corresponding to its complete class signature rather than receiving separate per-instance quotas.

\paragraph{Generated-count scaling.}
For a selector field $a$, let $n_a^{\mathrm{gen}}(v)$ be the raw count of generated images with value $v$. The effective count used in the frequency penalty is $\tilde{n}_a^{\mathrm{gen}}(v)=\gamma_r n_a^{\mathrm{gen}}(v)$. The main ImageNet10 and KITTI repair runs use $\gamma_r=1$, while the CelebA and COCO repair runs use $\gamma_r=1/r$, where $r$ is the retained source ratio. This scaling is applied before the source and generated counts are combined in Eq.~\ref{eq:frequency-penalty}.

\paragraph{Generation and assignment cost.}
On one NVIDIA H200 GPU, generating 1,000 T2I repair images with ControlNet takes approximately 20--30 minutes, and assigning repair-relevant attributes to these 1,000 generated images with Qwen3-VL takes approximately 5--10 minutes. The assignment time depends on the number of queried attributes, as reflected by the normalized source-pool measurements below. These shared upstream costs are reported separately from source selection and repair training because every selector uses the same fixed generated set. The source-pool assignment times reported below use the same one-H200 setting.

\subsection{End-to-End Cost Accounting}
\label{app:end-to-end-cost}

Table~\ref{tab:selection-efficiency} isolates source-subset selection after attribute records are available. We additionally measure the one-time source-pool assignment cost and the complete first repair run. T2I generation and generated-image assignment are reported separately because every selector uses the same fixed generated set. Source-pool assignment is charged to \model because it is a method-specific prerequisite when suitable metadata do not already exist.

\begin{table*}[t]
\centering
\caption{Measured source-pool attribute-assignment cost on one NVIDIA H200 GPU. The normalized time divides the total assignment time by the source-pool size in thousands and the number of assigned attributes.}
\label{tab:source-assignment-cost}
\small
\setlength{\tabcolsep}{4pt}
\begin{tabular}{l r r r r}
\toprule
\rowcolor{sectiongray} \textbf{Dataset} & \textbf{Source images} & \textbf{Attributes} & \textbf{Assignment time} & \textbf{min / 1K / attr.} \\
\midrule
KITTI & 5,000 & 6 & 32.77 min & 1.09 \\
ImageNet10 & 10,400 & 3 & 29.35 min & 0.94 \\
COCO & 20,000 & 5 & 90.00 min & 0.90 \\
\bottomrule
\end{tabular}
\end{table*}

The first-run totals in Tables~\ref{tab:kitti-end-to-end-cost} and~\ref{tab:coco-end-to-end-cost} include the applicable source-side preprocessing, subset selection, repair training, and evaluation. Every \model entry conservatively includes a fresh assignment of the full source pool. These tables give upper bounds for repeated repair cycles on the same task-level vocabulary.

\begin{table*}[t]
\centering
\caption{KITTI first-repair wall-clock time. Every \model row includes 32.77 minutes for fresh assignment of all 5,000 source images.}
\label{tab:kitti-end-to-end-cost}
\footnotesize
\setlength{\tabcolsep}{2.5pt}
\resizebox{\linewidth}{!}{%
\begin{tabular}{c c c c c c c c >{\columncolor{myblue!8}}c}
\toprule
\rowcolor{sectiongray} \textbf{Ratio} & \textbf{Random} & \textbf{Stratified} & \textbf{Entropy} & \textbf{CoreSet} & \textbf{BADGE} & \textbf{UHerding} & \textbf{PPAL} & \cellcolor{myblue!8}\textbf{Ours} \\
\midrule
0.1 & 42.17 min & 42.18 min & 48.82 min & 1.37 h & 1.12 h & 43.47 min & 43.53 min & 1.01 h \\
0.3 & 43.84 min & 43.85 min & 50.28 min & 1.44 h & 1.20 h & 45.08 min & 45.63 min & 1.09 h \\
0.5 & 45.18 min & 45.18 min & 51.73 min & 1.49 h & 1.27 h & 46.42 min & 47.47 min & 1.11 h \\
0.7 & 47.04 min & 47.05 min & 53.46 min & 1.56 h & 1.35 h & 48.16 min & 49.73 min & 1.13 h \\
0.9 & 48.86 min & 48.87 min & 55.53 min & 1.63 h & 1.42 h & 50.34 min & 51.46 min & \textbf{1.15 h} \\
\bottomrule
\end{tabular}%
}
\end{table*}

\begin{table*}[t]
\centering
\caption{COCO first-repair wall-clock time. Every \model row includes 90 minutes for fresh assignment of all 20,000 source images.}
\label{tab:coco-end-to-end-cost}
\footnotesize
\setlength{\tabcolsep}{2.5pt}
\resizebox{\linewidth}{!}{%
\begin{tabular}{c c c c c c c c >{\columncolor{myblue!8}}c}
\toprule
\rowcolor{sectiongray} \textbf{Ratio} & \textbf{Random} & \textbf{Stratified} & \textbf{Entropy} & \textbf{CoreSet} & \textbf{BADGE} & \textbf{UHerding} & \textbf{PPAL} & \cellcolor{myblue!8}\textbf{Ours} \\
\midrule
0.1 & 4.41 min & 4.42 min & 10.98 min & 16.18 min & 19.63 min & 5.08 min & 6.97 min & 1.57 h \\
0.3 & 8.18 min & 8.18 min & 14.74 min & 39.86 min & 49.76 min & 8.93 min & 17.86 min & 1.62 h \\
0.5 & 12.37 min & 12.36 min & 19.04 min & 1.08 h & 1.56 h & 13.08 min & 29.47 min & 1.66 h \\
0.7 & 15.73 min & 15.75 min & 22.27 min & 1.45 h & 2.29 h & 16.34 min & 37.63 min & 1.69 h \\
0.9 & 19.66 min & 19.65 min & 26.28 min & 1.87 h & 3.01 h & 20.18 min & 40.94 min & \textbf{1.72 h} \\
\bottomrule
\end{tabular}%
}
\end{table*}

Lightweight selectors remain cheaper at small ratios. At ratio $0.9$, however, \model is both faster and more accurate than CoreSet and BADGE on COCO and KITTI, even after full-pool assignment is charged. For example, the COCO times are 1.72 hours for \model, 1.87 hours for CoreSet, and 3.01 hours for BADGE; their test AP values are $0.6355$, $0.6202$, and $0.6213$, respectively. On KITTI, the corresponding times are 1.15, 1.63, and 1.42 hours, with mAP50--95 values $0.5850$, $0.5781$, and $0.5810$.

For a fixed task-level attribute vocabulary, source-pool assignment is a one-time preprocessing step. The resulting metadata can be reused across source ratios, attribute-aware selectors, and later repair iterations on the same task. The full source pool therefore does not need to be annotated again for every repair. Tables~\ref{tab:kitti-reused-cost} and~\ref{tab:coco-reused-cost} report the resulting wall-clock time after reusing source-attribute metadata.

\begin{table*}[t]
\centering
\caption{KITTI wall-clock time after reusing source-attribute metadata.}
\label{tab:kitti-reused-cost}
\footnotesize
\setlength{\tabcolsep}{2.5pt}
\resizebox{\linewidth}{!}{%
\begin{tabular}{c c c c c c c c >{\columncolor{myblue!8}}c}
\toprule
\rowcolor{sectiongray} \textbf{Ratio} & \textbf{Random} & \textbf{Stratified} & \textbf{Entropy} & \textbf{CoreSet} & \textbf{BADGE} & \textbf{UHerding} & \textbf{PPAL} & \cellcolor{myblue!8}\textbf{Ours} \\
\midrule
0.1 & 42.17 min & 42.18 min & 48.82 min & 1.37 h & 1.12 h & 43.47 min & 43.53 min & 42.20 min \\
0.3 & 43.84 min & 43.85 min & 50.28 min & 1.44 h & 1.20 h & 45.08 min & 45.63 min & 43.87 min \\
0.5 & 45.18 min & 45.18 min & 51.73 min & 1.49 h & 1.27 h & 46.42 min & 47.47 min & 45.20 min \\
0.7 & 47.04 min & 47.05 min & 53.46 min & 1.56 h & 1.35 h & 48.16 min & 49.73 min & 47.07 min \\
0.9 & 48.86 min & 48.87 min & 55.53 min & 1.63 h & 1.42 h & 50.34 min & 51.46 min & 48.89 min \\
\bottomrule
\end{tabular}%
}
\end{table*}

\begin{table*}[t]
\centering
\caption{COCO wall-clock time after reusing source-attribute metadata.}
\label{tab:coco-reused-cost}
\footnotesize
\setlength{\tabcolsep}{2.5pt}
\resizebox{\linewidth}{!}{%
\begin{tabular}{c c c c c c c c >{\columncolor{myblue!8}}c}
\toprule
\rowcolor{sectiongray} \textbf{Ratio} & \textbf{Random} & \textbf{Stratified} & \textbf{Entropy} & \textbf{CoreSet} & \textbf{BADGE} & \textbf{UHerding} & \textbf{PPAL} & \cellcolor{myblue!8}\textbf{Ours} \\
\midrule
0.1 & 4.41 min & 4.42 min & 10.98 min & 16.18 min & 19.63 min & 5.08 min & 6.97 min & 4.44 min \\
0.3 & 8.18 min & 8.18 min & 14.74 min & 39.86 min & 49.76 min & 8.93 min & 17.86 min & 8.20 min \\
0.5 & 12.37 min & 12.36 min & 19.04 min & 1.08 h & 1.56 h & 13.08 min & 29.47 min & 12.39 min \\
0.7 & 15.73 min & 15.75 min & 22.27 min & 1.45 h & 2.29 h & 16.34 min & 37.63 min & 15.77 min \\
0.9 & 19.66 min & 19.65 min & 26.28 min & 1.87 h & 3.01 h & 20.18 min & 40.94 min & 19.68 min \\
\bottomrule
\end{tabular}%
}
\end{table*}

With metadata reuse, \model has nearly the same wall-clock cost as Random and Stratified at every ratio, while remaining substantially faster than CoreSet and BADGE. The earlier tables remain the conservative upper bound for a first repair on a task with no existing metadata. Public metadata can remove this cost entirely when it covers the required vocabulary; CelebA, for example, provides 40 binary attributes with the dataset. Therefore, the $39$--$283\times$ values in Table~\ref{tab:selection-efficiency} describe source-subset selection speed after attribute records are available, not end-to-end pipeline speed.
